% Separate amendment: depends on the full-CPsh version of the polar theorem.
% Do not insert against the older, first-stage-only theorem.
\begin{corollary}
\label{cor:ch16-cpsh-component-curvature}
Let $K$ be any complete non-Archimedean field, possibly trivially valued,
and let $X$ be a good separated usual-strict boundaryless pure
$n$-dimensional $K$-analytic space, with $n\geq1$.
For an arbitrary ordinary open $W\subseteq X$, let $(W_i)_i$ be the
reduced analytic irreducible components of $W$ over $K$, let
$j_i:W_i\hookrightarrow W$, and set
$\delta_i=(j_i)_*\delta_{W_i}$.
Then, for every $u,v\in\operatorname{CPsh}_W(W)$,
\begin{equation}
 \dprimeddprime(u\delta_W)=\dprimeddprime(v\delta_W)
 \quad\Longleftrightarrow\quad
 \dprimeddprime(u\delta_i)=\dprimeddprime(v\delta_i)
 \quad\text{for every }i .
 \label{eq:ch16-cpsh-component-curvature}
\end{equation}
Here $\operatorname{CPsh}$ is the full continuous uniform-closure sheaf,
not just the class $\mathscr A$ of
\Cref{sec:ch14-component-curvature}.
Neither reducedness nor normality is required, and component
intersections may have dimension $n-1$.
\end{corollary}

\begin{proof}
We first spell out the ordinary-local use of the full-$\operatorname{CPsh}$
analytic-nullity conclusion of
\Cref{cor:ch16-cpsh-analytic-nullity}.
At a point of $W$, the barrier construction in that proof can be
performed inside the prescribed ordinary neighbourhood $W$:
choose an affinoid neighbourhood there, then a smaller rational
affinoid neighbourhood, and use the regularized maximum to obtain a
smooth psh function with a compact positive closed sublevel.
The potentials are needed only on a neighbourhood of this compact
sublevel, which lies in $W$.

If $Z$ is closed analytic in $W$ and has local dimension less than $n$,
local prime avoidance gives an analytic function $f$ vanishing on $Z$
but on no top-dimensional local component.  Its logarithm is locally
integrable for canonical top-form measures by the local top-form
support property used in the same analytic-nullity proof.
The smooth convex truncations of $\log|f|$ are defined on an ordinary
neighbourhood inside $W$ and decrease to $\log|f|$.
The ordinary-open weighted polar assertion therefore applies with
these truncations and with the given $\operatorname{CPsh}$ potentials,
without extending any of them beyond $W$.
It gives zero total-variation mass on $Z$ for each curvature coefficient.
A finite smooth partition near a compact test support gives the
assertion throughout $W$.
This is use of the local proof on an ordinary subopen; it is not an
assertion that $W$ remains usual-strict over a trivially valued field.

The analytic component geometry and canonical normalization used in
\Cref{prop:ch14-component-curvature} apply locally on $W$.
Write $m_i>0$ for its generic multiplicities.  The component family is
locally finite near compact sets and
\[
                 \delta_W=\sum_i m_i\delta_i .
\]
The sets
\[
 B_i=W_i\cap\bigcup_{j\ne i}W_j,
 \qquad B=\bigcup_iB_i
\]
are closed analytic in $W$ of local dimension less than $n$.
These are the actual components of $W$; they need not be restrictions
of a fixed global component family.

The reverse implication in
\eqref{eq:ch16-cpsh-component-curvature} follows by multiplying the
component equalities by $m_i$ and summing locally finitely.
For the forward implication put
\[
 S=\dprimeddprime(u\delta_W)=\dprimeddprime(v\delta_W),
 \qquad S_i^u=\dprimeddprime(u\delta_i),\qquad
        S_i^v=\dprimeddprime(v\delta_i).
\]
The full-sheaf positivity argument in the proof of
\cref{itm:ch14-component-curvature:common} gives positive separately
closed currents $S_i^u,S_i^v$ on $W$, and
\[
                    S=\sum_i m_iS_i^u=\sum_i m_iS_i^v .
\]
For each compactly supported positive smooth
$(n-1,n-1)$-form $\alpha$, these positive decompositions give the
Radon-measure inequalities
\[
 0\leq S_i^u\wedge\alpha\leq m_i^{-1}S\wedge\alpha,
 \qquad
 0\leq S_i^v\wedge\alpha\leq m_i^{-1}S\wedge\alpha.
\]
The local analytic-nullity assertion just proved, with the one factor
$u$ and $Z=B$, says that $S\wedge\alpha$ has zero mass on $B$.
Thus both component measures have zero mass on $B_i$.
The symmetric part of an arbitrary compact smooth test is locally a
difference of two positive tests, and the antisymmetric part pairs to
zero.  Consequently every signed coefficient of
$R_i=S_i^u-S_i^v$ has total variation zero on $B_i$.

Outside $W_i$, the current $R_i$ vanishes by support.  On the ordinary
open $W\setminus\bigcup_{j\ne i}W_j$, the equality of canonical
curvatures reads $m_iR_i=0$.  These two opens cover $W\setminus B_i$,
so $\Supp R_i\subseteq B_i$.
Every coefficient measure of $R_i$ is therefore supported on $B_i$,
by smooth-test density, and has zero total variation there by the
preceding paragraph.  It is zero.
A compact smooth scalar cutoff equal to one near each test support
recovers $\langle R_i,\alpha\rangle=0$, so $R_i=0$ for every $i$.
This proves the forward implication.
\end{proof}

The corollary compares the specified potential curvatures on the
original analytic space, including its nonreduced multiplicities.
It does not classify arbitrary positive closed currents.
No scalar-current classification, extension of a current from $W$,
normalization pullback, or scalar extension of the ground field is used.
In particular it does not by itself settle the all-current,
all-bidegree question in \Cref{prob:frontier-cross-fan-upstairs}.
